\documentclass[a4paper]{article}
\usepackage{amsfonts}
\usepackage{amsmath}
\usepackage{amssymb}
\usepackage{graphicx}     

\begin{document}
  
\noindent \large Auteur: Abdellah El Moussaoui \\
\noindent \large Studierichting: Informatica\\
\noindent \large Oefeningenreeks 3F nr. 1.11\\

\medskip

\normalsize

$f(x)=x+1+\dfrac{3}{x-7}$\\

$f(x)=0 \Leftrightarrow x+1=-\dfrac{3}{x-7}\Leftrightarrow(x+1)(x-7)=-3$

$\Leftrightarrow (x-3)^2=13 \Leftrightarrow x=3\pm\sqrt{13}$\\

$f'(x)=1-\dfrac{3}{(x-7)^2}$\\

VA:\\

$x=7$\\

$\lim\limits_{x\to7+} f(x) = +\infty$\\

$\lim\limits_{x\to7-} f(x) = -\infty$\\

HA:\\

$\lim\limits_{x\to\infty} f(x) = +\infty$\\

Geen HA!\\

SA:\\

$\lim\limits_{x\to\infty} \dfrac{f(x)}{x}$

$= \lim\limits_{x\to\infty} \dfrac{x+1+\dfrac{3}{x-7}}{x}$\\

$=\lim\limits_{x\to\infty} 1 + \dfrac{1}{x} + \dfrac{3}{x(x-7)}$\\

$ = 1$\\

$\lim\limits_{x\to\infty} f(x)-x = 1$\\

Dus $f(x)$ heeft een schuine asymptoot, namelijk $y=x+1$\\

$f(x)-(x+1) = 0 \Leftrightarrow x\in \emptyset$\\

$\begin{array}{r|ccccccc}
           & -\infty &          & 7     &       & +\infty \\
    \hline  
    f(x)-(x+1)  &   -    & - & |& +   &  +   \\
\end{array}$ \\

Dus $f(x)$ is onder $y=x+1$ voor $x \in ]-\infty,7[$ en boven $y=x+1$ voor $x\in ]7,+\infty[ .$\\

\begin{thebibliography}{999}
%\bibitem{a} Referentie
\end{thebibliography}
\end{document} 